\section{Introduction}

The question of what happens when an autonomous multi-agent system encounters a condition outside its design envelope has transitioned from theoretical to operational. Autonomous agent populations are now deployed in agricultural resource management, clinical decision support, financial clearing, and critical logistics --- domains where the cost of an uncontrolled failure is measured in lives, capital, and institutional trust. As agent populations grow in recursive depth and operational scope, the absence of a structural mechanism compelling human involvement for unanticipated exception conditions has become the central vulnerability of agentic infrastructure.

We introduce the \textbf{Bailout Protocol}: the mandatory, architecturally enforced escalation mechanism of the \bxthree{} Framework. It compels every unresolved exception state in any \bxthree{} node to propagate upward to a named human principal, bypassing all intermediate machine actors, with full forensic attribution at every step. It is not a feature added to an existing architecture. It is the mechanism by which the \bxthree{} Framework makes its upstream accountability guarantee structurally operative at runtime.

\subsection{Why Autonomous Systems Require a Mandatory Bail-Out Mechanism}

Autonomous multi-agent systems operate within provisioned operating ranges $R(n)$ defined at node provisioning. Within $R(n)$, bounded reasoning and constrained execution are sufficient to maintain safe, productive operation. Outside $R(n)$, the system encounters conditions it was not designed to resolve --- unanticipated input distributions, emergent behaviors from recursive agent spawning, cross-domain inputs that no individual node's training data anticipated, or cascading failures from upstream components. These are not exceptional in the statistical sense; they are inevitable in the operational sense. Any system deployed at sufficient scale and recursive depth will encounter them.

The standard response in current architectures is discretionary escalation: a system \emph{may} flag an anomaly, \emph{may} notify an operator, and \emph{may} pause execution pending review. This model fails for three structural reasons. First, it makes human oversight contingent on machine actors' willingness to recognize that they have exceeded their competence --- a condition that is not reliably detectable by the actors themselves. Second, it permits absorption at intermediate tiers: a supervisory machine actor may absorb the exception, attempt resolution, and either succeed silently or fail in a manner that is never propagated upward. Third, it provides no forensic record of what occurred during escalation, making post-hoc accountability impossible when the escalation pathway was incomplete or suppressed.

These are not implementation failures amenable to better engineering. They are architectural failures: the accountability relationship between machine and human is defined by discretion rather than compulsion, and discretion can always select against escalation when escalation is costly, inconvenient, or reputationally threatening.

The Bailout Protocol replaces discretion with compulsion. A node enters the protocol not by choosing to escalate, but by observing that its current state exceeds $R(n)$ and that it cannot produce a \fact-layer-approved resolution within a bounded timeout. Entry to the protocol is non-discretionary and atomic. The exception state is then propagated upward through the accountability chain via a parent-chain bypass that contains no machine-actor terminus. The human principal $h(n)$ receives the exception, issues a binding resolution directive, and the protocol closes with a full forensic record.

\subsection{What Happens When Systems Operate Without an Architectural Fallback}

The consequence of operating without a mandatory bail-out mechanism is not incremental degradation. It is a category shift: from accountable autonomous operation to unconstrained autonomous collapse. We identify three failure modes that are structurally prevented by the Bailout Protocol but are individually sufficient to produce uncontrolled outcomes in systems lacking it.

The first is \textbf{exception absorption at intermediate tiers}. A supervisory machine actor receives an exception, evaluates it against its own operating range, and --- finding the exception within its apparent scope --- attempts resolution without escalation. The resolution is incorrect, or it addresses a symptom rather than a cause, and the exception state propagates laterally into sibling subsystems while the absorbing node reports nominal operation. The human principal receives no signal. The failure compounds silently.

The second is \textbf{delayed systemic failure through recursive depth}. A deeply spawned agent population encounters a cross-domain input that none of its constituent nodes were provisioned to handle. Each node individually determines that the input is marginally within its $R(n)$ interpretation, so no individual bailout trigger fires. The exception state manifests as an emergent property of the population --- a pattern visible at the system level but invisible at any single node. No escalation occurs because the trigger condition never fires at the node level. The system drifts toward a critical threshold and collapses without human awareness.

The third is \textbf{suppression under operational pressure}. A node correctly identifies an exception condition requiring escalation but faces operational pressure --- a deadline, a performance metric, or an explicit directive from a supervisory node --- not to escalate. In a discretionary escalation model, the node may suppress the escalation and attempt to proceed. If the exception is consequential, the suppression produces a failure that is undetectable at the system level because the escalation that would have revealed it was never initiated.

Each of these failure modes is structurally distinct, but they share a common root: the accountability chain between machine and human is defined by machine actors' self-assessment rather than by architectural constraint. The Bailout Protocol eliminates all three by making the chain a structural property of the system rather than a behavioral property of its components.

\subsection{Accountability as an Architectural Fallback}

The core insight motivating the Bailout Protocol is that accountability must have an architectural fallback --- not a discretionary one. In the \bxthree{} Framework, accountability is not a governance policy that system designers implement at their discretion. It is a structural property of the three-layer architecture: every node has a \purpose{} layer carrying intent and final authority, a \bounds{} layer carrying bounded reasoning and constrained execution, and a \fact{} layer carrying deterministic enforcement and forensic auditability.

When the \bounds{} layer encounters a condition outside $R(n)$ it cannot resolve, the functional separation itself dictates that the decision must return to the \purpose{} layer. There is no other layer with the authority to adjudicate it. The Bailout Protocol is not a feature added to this separation; it is the mechanism by which the separation enforces its own accountability guarantee. The human anchor $h(n)$ is not a design choice --- it is a structural necessity.

This distinguishes \bxthree{} from prior escalation architectures. The Tiered Agentic Oversight model~\cite{kim2025tao} permits absorption at intermediate machine tiers, making human oversight a high-tier option rather than a compulsory terminus. The Bailout Protocol forbids this absorption by architectural constraint: the distinction is the difference between a system that \emph{may} consult a human and a system that \emph{cannot proceed} without one.

The upstream accountability guarantee of the \bxthree{} Framework --- that \bxthree{} systems fail upward into human consciousness, never downward into autonomous chaos --- requires the Bailout Protocol to be structurally operative. Without it, the guarantee is a claim without a mechanism. With it, the guarantee is enforceable for every unanticipated exception condition across the full recursive depth of the agent population.

\subsection{Paper Roadmap}

The remainder of this paper is organized as follows. Section~\ref{sec:bailout} provides the complete formal specification of the Bailout Protocol: its deterministic state machine $\mathcal{B}$, the trigger condition governing entry to the protocol, the escalation algorithm, the parent-chain bypass enforcing the no-machine-actor property, and the timeout handling guaranteeing liveness under all failure conditions. Section~\ref{sec:implementation} addresses the protocol's realization within the \bxthree{} three-layer architecture, showing how the \purpose{}, \bounds{}, and \fact{} layers each participate in protocol execution and how the Forensic Ledger is maintained as an immutable evidentiary record. Section~\ref{sec:properties} formally states the Safety, Liveness, and Accountability correctness properties, derives them from the protocol's specification, and provides proof sketches under the stated structural assumptions. Section~\ref{sec:related} surveys related work across autonomous agent architectures, escalation mechanisms, and accountability frameworks. Section~\ref{sec:limitations} concludes with structural limitations and a directed research agenda.
